\begin{definition}
A \emph{quasiregular-theory package} is a datum consisting of proofs of the
following four statements.
\begin{enumerate}
\item[(a)] (Continuity.)  For every $K$, every open $\Omega\subset\mathbb{C}$
and every $f$ which is $K$-quasiregular on $\Omega$ in the sense of
\noderef{QuasiregularOn}, there is $g\colon\mathbb{C}\to\mathbb{C}$ with
$g=f$ almost everywhere on $\Omega$ and $g$ continuous on $\Omega$.

\item[(b)] (Discreteness of the zero set.)  For every $K$, every open
connected $\Omega$, every $f$ which is $K$-quasiregular on $\Omega$, and every
$g$ with $g=f$ almost everywhere on $\Omega$ and $g$ continuous on $\Omega$
such that $g$ does not vanish identically on $\Omega$, every point
$z_0\in\Omega$ has a punctured neighbourhood in $\Omega$ on which $g$ does not
vanish: there is $r>0$ such that $g(z)\neq0$ for all $z\in\Omega$ with
$z\neq z_0$ and $|z-z_0|<r$.

\item[(c)] (Factorization near a zero.)  For every $K$, every open connected
$\Omega$, every $f$ which is $K$-quasiregular on $\Omega$, every $g$ with
$g=f$ almost everywhere on $\Omega$ and $g$ continuous on $\Omega$ such that
$g$ does not vanish identically on $\Omega$, and every $z_0\in\Omega$ with
$g(z_0)=0$, there are an integer $n\ge1$, a constant $K'$, an open set $V$
with $z_0\in V\subset\Omega$ and a map $\chi$ which is a $K'$-quasiconformal
homeomorphism of $V$ in the sense of \noderef{QuasiconformalOn}, such that
\[
  \chi(z_0)=0\qquad\text{and}\qquad g(z)=\bigl[\chi(z)\bigr]^n\ \text{ for }z\in V .
\]

\item[(d)] (The inverse of a quasiconformal homeomorphism.)  For every $K$,
every open $V$, every $\chi$ which is a $K$-quasiconformal homeomorphism of
$V$, and every $H$ with $H(\chi(z))=z$ for all $z\in V$, the map $H$ belongs
to $W^{1,2}_{\mathrm{loc}}(\chi(V))$ in the sense of \noderef{MemWLocOn}, is
real-differentiable at almost every point of $\chi(V)$, and satisfies
$J(\xi)>0$ for almost every $\xi\in\chi(V)$, where $J$ is the Jacobian of
\noderef{Jac}.
\end{enumerate}

This is item (C2) of the paper's note for the formalization, quoted there from
[Ah] and [LV]: ``A $K$-quasiregular mapping is continuous; if it is
not identically zero its zero set is discrete.  Near a zero $z_0$ it factors
as $f=\chi^n$ with $n\ge1$ and $\chi$ a quasiconformal homeomorphism of a
neighborhood of $z_0$ with $\chi(z_0)=0$; the inverse $H$ of a quasiconformal
homeomorphism belongs to $W^{1,2}_{\mathrm{loc}}$, has Jacobian $J>0$ almost
everywhere, and satisfies (17).''  Clause (c) is the factorization (15) of
Section~3.  The note directs that such quoted results be carried as explicit
hypotheses rather than axioms; this definition is that bundle, so that a
theorem needing (C2) takes a quasiregular-theory package as an argument.

Three points of reading.

First, why the first three clauses speak of a representative $g$ rather than
of $f$.  Quasiregularity as \noderef{QuasiregularOn} defines it is a property
of the almost-everywhere class of $f$ and of nothing more.  Its two clauses
are membership in $W^{1,2}_{\mathrm{loc}}$, which by \noderef{MemWLocOn} is a
family of integral identities and indicator $L^2$-conditions, and an
almost-everywhere inequality between the weak Wirtinger derivatives of $f$,
which by \noderef{IsWeakWirtingerOn} depend only on that class as well.  Both
are unchanged when $f$ is altered on a set of measure zero, so neither can
force $f$ to be continuous at any point.  The function $f$ equal to $1$ at the
origin and to $0$ elsewhere is $K$-quasiregular on $\mathbb{C}$ for $K=1$: it
vanishes almost everywhere, so the integral identities and the indicator
conditions are those of the zero function, and the zero function is a
distributional $\partial_0$- and $\partial_1$-derivative of $f$ on
$\mathbb{C}$, since $\int f\,\partial_j\phi\,d\sigma=0=-\int0\cdot\phi\,
d\sigma$ for every test function; so both weak Wirtinger derivatives of $f$
may be taken to be $0$ and the inequality reads $0\le0$.  Yet $f$ is
discontinuous at the origin, has $f(1)=0$ with $f$ vanishing on a whole
punctured neighbourhood of $1$, and cannot factor as $\chi^n$ near $1$ with
$\chi$ injective and continuous.  So the three clauses stated of $f$ would be
false, the package would be uninhabited, and every theorem taking one would be
vacuously true; `paper.tex:107` forbids exactly that, requiring the quoted
results not to be vacuous.  Stated of a continuous representative they are the
literature's statements: a quasiregular mapping \emph{is} a
$W^{1,2}_{\mathrm{loc}}$ map with a distortion bound on its weak derivatives,
and ``it is continuous'' means that it has a continuous representative.  The
representative is unique: two continuous functions on an open set agreeing
almost everywhere agree everywhere on it.

It matters here that the distortion inequality of \noderef{QuasiregularOn} is
read for the weak derivatives.  Were it read for the pointwise Wirtinger
derivatives of the given representative, which are built from
$\mathrm{d}f$ and vanish wherever $f$ is not differentiable, the hypothesis
$K$-quasiregular would be met by representatives with no continuous
representative at all --- for instance by $f_0(z)=\log\log(1/|z|)$ on
$\{|z|<1/2\}$, which lies in $W^{1,2}$ there but is essentially unbounded at
the origin, altered by $1$ on the countable dense null set
$\mathbb{Q}+i\mathbb{Q}$, so that it becomes nowhere continuous and hence
nowhere differentiable while both pointwise derivatives vanish identically ---
and clause (a) would then be false, making this bundle uninhabited.  With the
weak reading, $(a)$ asserts no more and no less than the literature's
continuity theorem.

Second, the zero set in (b) and (c) is the representative's.  The set
$f^{-1}(0)$ is not determined by $f$ modulo null sets, whereas $g^{-1}(0)$ is;
and the paper itself writes the exceptional set as
$f^{-1}(0)=\{z\in\Omega:\nabla u(z)=0\}$ at paper.tex:156, that is, as the zero
set of the continuous complex gradient.  Only clause (a) produces a
representative, so a consumer of the package introduces a single $g$ and runs
(b), (c) and the bundle (C3) on that same $g$.

Third, clause (d) needs no such reading, and the final clause ``and satisfies
(17)'' of the quoted (C2) is deliberately \emph{not} among the assumptions
above.  The hypotheses of (d) constrain $\chi$ through
\noderef{QuasiconformalOn}, which carries genuine pointwise content ---
injectivity, continuity, openness of the image and a continuous inverse ---
and its three conclusions are themselves almost-everywhere statements, so it
is inhabited as it stands.  Equation (17) is derived in the tablet from the
chain rule for the Wirtinger operators, wherever $H$ and $\chi$ are
differentiable at matched points and $J\neq0$, and the almost-everywhere
differentiability and positivity of $J$ which make that derivation applicable
almost everywhere are assumed in (d).  Omitting a clause that is proved rather
than quoted leaves the package no stronger than the literature statement, which
is what the note requires.
\end{definition}
